\section{Periodic Orbit Design}
\label{sec:orbit}

\subsection{An Overview of Orbit Families}

Usable cislunar orbits fall into four layers by distance from the Moon, with
scale as the main thread: the farther from the Moon, the larger Earth's share
of the gravity field, the dynamics transition from two-body Keplerian to
three-body, and orbits turn from stable to unstable
(Table~\ref{tab:families}). e2m2e provides a unified family-generation entry,
\texttt{orbit\_family\_generation()}, covering eight orbit families; Lyapunov
and resonant orbits (RO) are additionally supported at the differential
correction strategy layer.

\begin{table}[htbp]
\centering\small
\caption{Periodic orbit families}
\label{tab:families}
\begin{tabularx}{\textwidth}{@{}llX@{}}
\toprule
\textbf{Family} & \textbf{Location} & \textbf{Characteristics and uses} \\
\midrule
DRO & Moon-centered & Distant retrograde orbit; neutrally stable with low
maintenance cost; candidate for lunar orbital stations \\
DPO & near L2 & Distant prograde orbit; unstable; usable for far-side
communication relay \\
Halo / NRHO & L1/L2 & Three-dimensional periodic orbits; NRHOs are the
high-amplitude near-rectilinear form ($A_z > 0.1$ dimensionless), adopted by
Gateway \\
Axial & L1/L2 & G\'omez Type B bifurcation; axial motion, coverage
complementary to Halo \\
Lyapunov & L1/L2/L3 & Planar periodic orbits \\
Lissajous & L1/L2/L3 & Quasi-periodic (incommensurate in-plane/out-of-plane
frequencies), generated on the nonlinear center-reduced flow \\
SPO / LPO & L4/L5 & Short-/long-period orbits; large-amplitude LPOs take a
horseshoe shape \\
Horseshoe & L4/L5 & Large-amplitude members of the LPO chain \\
RO & --- & Resonant orbits, departing from the $y$-axis symmetry \\
\bottomrule
\end{tabularx}
\end{table}

Periodicity semantics are strictly distinguished: Halo, NRHO, Axial, SPO,
LPO, Horseshoe, and DRO members are strictly periodic; Lissajous members are
quasi-periodic bounded trajectories and are not forced into periodic closure.

\subsection{Initial Guesses: Starting from Analytic Approximations}

The first building block of periodic orbit design is the initial guess, for
which e2m2e has three sources:

\begin{itemize}
  \item \textbf{Richardson third-order analytic approximation} (1980): a
        third-order analytic solution for Halo orbits near the collinear
        libration points. The internal pipeline: solve the libration-point
        distance parameter $\gamma$ (quintic equation, Brent's method) $\to$
        compute the in-plane frequency $\omega_p$ $\to$ assemble the
        third-order coefficients and solution. Small-amplitude seeds (e.g.
        dimensionless $z_0 = 0.001$) are the most accurate and are then
        amplified by continuation;
  \item \textbf{Seed orbits}: the DRO family starts from a standard seed
        (amplitude $\approx 90786$ km), differentially corrected with the
        $x$-axis crossing fixed and continued bidirectionally, covering an
        envelope of 1737--110000 km;
  \item \textbf{Normal-form characteristic parameters}: the dynamics near a
        libration point are reduced to action--angle characteristic
        parameters $[q_1, p_1, I_2, \theta_2, I_3, \theta_3]$, describing
        with a few constants an orbit that would otherwise need six state
        components evolving in time. The pipeline: ephemeris initial state
        $\to$ dynamical substitute orbit $\to$ quasi-Floquet transformation
        $\to$ center-manifold reduction $\to$ characteristic parameters. The
        quasi-Floquet reduction offers a 36-dimensional matrix method or a
        21-dimensional $\mathfrak{sp}(6)$ Lie-algebra parameterization
        (symplectic by construction); the center-manifold Lie transformation
        is truncated at order 10 by default.
\end{itemize}

\subsection{Differential Correction: STM Newton Iteration via Symmetries}

Differential correction iteratively adjusts the initial conditions so that
the orbit satisfies periodicity constraints, at its core an STM-based Newton
iteration. e2m2e exploits the three mirror symmetries of the CR3BP to shrink
the problem---for example, XZ-plane symmetry means that if
$(x, y, z, \dot x, \dot y, \dot z)$ is a solution, so is
$(x, -y, z, -\dot x, \dot y, -\dot z)$; a Halo orbit therefore only needs to
depart perpendicularly from the XZ plane ($y = 0$) and be corrected to cross
it perpendicularly again half a period later.

Different combinations of symmetry and fixed parameters form 11 correction
strategies (Lyapunov/DRO with fixed $x_0$, fixed-period orbits with fixed
$T_{1/2}$, RO with fixed $y_0$, Halo with fixed $z_0$ or $x_0$, Axial with
fixed $v_{z0}$, SPO/LPO at L4/L5 without symmetry and full-period closure,
etc.). Architecturally, the strategy layer is strictly separated from the
solver: strategy functions only produce immutable configuration objects, and
execution sinks to the Rust CR3BP correction kernel.

Convergence failures are diagnosable: the history records per-step errors and
the termination cause (converged / correction below machine precision /
diverged / singular Jacobian / stagnated / converged to a parasitic root with
$T \approx 0$); a singular Jacobian indicates linearly dependent constraints
and calls for a different strategy.

\subsection{Multiple Shooting: Ephemeris Correction of Long Arcs}

Differential correction is sensitive to the initial guess, increasingly so
for longer arcs. Multiple shooting splits the orbit into $N$ independently
integrated segments and imposes continuity constraints at the patch points:
\begin{equation}
  \vect{x}_i(t_{i+1}) = \vect{x}_{i+1}(t_{i+1}),
\end{equation}
adjusting all segment initial states (optionally including time nodes)
simultaneously---turning ``long-arc sensitivity'' into ``many small
corrections''. The standard implementation reaches a tolerance of $10^{-10}$.

Patch-point correction under the ephemeris model uniformly goes through the
Rust multiple-shooting kernel, dispatched by orbit stability along two paths:

\begin{itemize}
  \item \textbf{Stable orbits} (DRO): a velocity-weighted two-level
        path---after correcting one revolution, free extrapolation remains
        bounded;
  \item \textbf{Unstable orbits} (Halo/NRHO/DPO): segmented full-arc
        shooting, integrating segment by segment to fill the ephemeris grid;
        NRHO fixes 1 revolution per segment, Halo at most 3, and DPO uses 64
        equally-spaced patch points per revolution.
\end{itemize}

Node sampling is likewise tailored to orbital physics: Halo uses
perilune-clustered sampling (5 nodes inside the perilune window plus 8
outside), because the STM is ill-conditioned where the velocity is large
near perilune; NRHO uses equal-time sampling.

\subsection{Continuation: From One Orbit to a Family}

Given a converged seed orbit, a family is generated by stepping along a
family parameter (Jacobi constant, amplitude, etc.), reconverging with
differential correction at each step. Two continuation schemes are provided:

\begin{itemize}
  \item \textbf{Natural-parameter continuation}: advances along a single
        parameter, each step seeded by the current solution;
  \item \textbf{Pseudo-arclength continuation (PAL)}: adds an arclength
        constraint and can pass turning points, suitable for families with
        non-monotonic parameter--amplitude relations (e.g. the Halo
        amplitude--period curve).
\end{itemize}

Failure handling is engineered: a non-converged step retries with a reduced
step size; if convergence still fails at the lower bound, that direction
terminates; only converged members are appended to the family, with
diagnostics reporting total/successful/failed steps. The L2 NRHO family uses
a dedicated ``fold then continue with fixed $x_0$'' arrangement; the Halo
family uses natural continuation with fixed $z_0$, positive and negative
$z_0$ giving the northern and southern families.

\subsection{Stability Analysis and Invariant Manifolds}

The stability of a periodic orbit is determined by its monodromy matrix (the
STM over one period):
\begin{equation}
  \mat{M} = \stm(T, 0).
\end{equation}
By the symplectic structure of Hamiltonian systems, eigenvalues appear in
reciprocal pairs $(\lambda, 1/\lambda)$. The Broucke stability index sums
each reciprocal pair:
\begin{equation}
  \nu = \lambda + \frac{1}{\lambda},
\end{equation}
with $|\nu| < 2$ stable and $|\nu| > 2$ unstable (the larger, the more
unstable). The relation of eigenvalues to $+1$, $-1$, and the unit circle
identifies bifurcation types (saddle-node / period-doubling / torus), with
family-level bifurcation detection supported. Numerical trustworthiness is
self-checked via symplectic residuals: if $|\det \mat{M} - 1|$ or
$\norm{\mat{M}^\top \mat{J} \mat{M} - \mat{J}}$ is significantly large, the
integration tolerance should be tightened and the analysis redone.

The real eigenvectors of the monodromy matrix also give the invariant
manifold directions: the stable manifold takes the $|\lambda| < 1$ direction
and the unstable manifold the $|\lambda| > 1$ direction. Transporting the
eigenvectors along the orbit by the STM and applying a $\pm\varepsilon$
perturbation (typically 50 km/DU) yields manifold seeds; integrating
backward/forward produces the manifold tubes used for low-energy transfer
stitching (Section~\ref{sec:transfer}).

\subsection{The Design Pipeline}

Chaining these algorithms together, the complete pipeline of one mission
orbit is:

\begin{center}
\begin{tikzpicture}[
  font=\footnotesize,
  stage/.style={draw=e2blue, thick, rounded corners=1.5mm, fill=e2light,
    minimum height=0.85cm, align=center, inner sep=4pt},
  arr/.style={-{Stealth[length=2mm]}, thick, e2blue},
]
\node[stage] (s1) {initial guess\\Richardson/seed};
\node[stage, right=0.24cm of s1] (s2) {differential\\correction};
\node[stage, right=0.24cm of s2] (s3) {continuation\\family};
\node[stage, right=0.24cm of s3] (s4) {phasing\\+ nodes};
\node[stage, right=0.24cm of s4] (s5) {batch frame\\conversion};
\node[stage, right=0.24cm of s5] (s6) {ephemeris\\correction};
\node[stage, right=0.24cm of s6] (s7) {high-fidelity\\propagation};
\draw[arr] (s1) -- (s2);
\draw[arr] (s2) -- (s3);
\draw[arr] (s3) -- (s4);
\draw[arr] (s4) -- (s5);
\draw[arr] (s5) -- (s6);
\draw[arr] (s6) -- (s7);
\end{tikzpicture}
\end{center}

All numerical heavy lifting (propagation, shooting, batch frame conversion,
batch ephemeris queries) runs on the Rust side, with Python only
orchestrating; the end-to-end generation of four typical orbits (LLO, ELFO,
DRO, Halo) takes under one minute in total with a release build.

\begin{keybox}{NominalOrbit: The Contract Between Design and Control}
A designed nominal orbit is delivered as a NominalOrbit structure: an
equally-spaced state table + a Floquet basis + projection factors.
Station-keeping algorithms consume this structure directly---the Floquet
analysis produced during design does not need to be recomputed during
control; the two stages are decoupled through a single data contract.
\end{keybox}
